% SEGMENT OLP-0036-S01
% Part:sets-functions-relations
% Chapter: size-of-sets
% Section: comparing-sizes

\documentclass[../../../include/open-logic-section]{subfiles}

\begin{document}

\olfileid{sfr}{siz}{car}

\olsection{வெவ்வேறு அளவுள்ள கணங்களும் காண்டோரின் தேற்றமும்}

\begin{explain}
இரண்டு கணங்கள் ஒரே அளவைக் கொண்டுள்ளன என்ற கருத்தைத்
துல்லியமாகக் கூறியுள்ளோம்.
ஒரு கணம் மற்றொன்றைவிடச் சிறியது என்ற கருத்தையும் துல்லியமாகக் கூறலாம்.
“சிறியது (அல்லது எண்ணொத்தது)” என்பதற்கான நமது வரையறையில்,
கணங்களுக்கு இடையிலான !!a{bijection} க்குப் பதிலாக,
முதல் கணத்திலிருந்து இரண்டாவது கணத்திற்குச் செல்லும்
!!a{injection} தேவைப்படும்.
இத்தகைய சார்பு இருந்தால், முதல் கணத்தின் அளவு இரண்டாவது கணத்தின்
அளவுக்குச் சமமாகவோ சிறியதாகவோ இருக்கும்.
உள்ளுணர்வின்படி, ஒரு கணத்திலிருந்து மற்றொன்றிற்குச் செல்லும்
!!a{injection}, சார்பின் வீச்சகத்தில் அதன் சார்பகத்தில் உள்ள அளவுக்குக்
குறைவில்லாத எண்ணிக்கையில் !!{element}s இருப்பதை உறுதிசெய்கிறது.
ஏனெனில் சார்பகத்தின் இரு !!{element}s வீச்சகத்தின் ஒரே
!!{element} க்கு அனுப்பப்படுவதில்லை.
\end{explain}

% SEGMENT OLP-0036-S02
\begin{defn}
!!a{injection} $f \colon A \to B$ இருக்கும்போது,
அப்போதும் அப்போது மட்டுமே, $A$ என்பது $B$ ஐ
\emph{விடப் பெரிதல்ல} எனப்படும்; இதனை $\cardle{A}{B}$ என எழுதுகிறோம்.
\end{defn}

% SEGMENT OLP-0036-S03
இது தற்சுட்டு, கடப்பு ஆகிய பண்புகளைக் கொண்ட தொடர்பு;
ஆனால் சமச்சீர்ப் பண்புடையதல்ல என்பது தெளிவு
(இது ஒரு பயிற்சியாக விடப்பட்டுள்ளது).
ஒரு கணம் மற்றொன்றைவிட (கண்டிப்பாகச்) சிறியது என்பதைக் கூறும்
ஒரு கருத்தையும் அறிமுகப்படுத்தலாம்.

\begin{defn}
!!a{injection} $f\colon A \to B$ இருந்தும்,
!!{bijection} $g\colon A
\to B$ எதுவும் இல்லாதபோது, அப்போதும் அப்போது மட்டுமே,
$A$ என்பது $B$ ஐ \emph{விடச் சிறியது} எனப்படும்.
இதனை $\cardless{A}{B}$ என எழுதுகிறோம்;
அதாவது $\cardle{A}{B}$ மற்றும் $\cardneq{A}{B}$.
\end{defn}

% SEGMENT OLP-0036-S04
இந்தத் தொடர்பு எங்கும் தற்சுட்டற்றது, கடப்புப் பண்புடையது என்பது தெளிவு.
(இது ஒரு பயிற்சியாக விடப்பட்டுள்ளது.)
இந்தக் குறியீட்டைப் பயன்படுத்தி,
$\cardle{A}{\Nat}$ என இருக்கும்போது, அப்போதும் அப்போது மட்டுமே,
ஒரு கணம் $A$ என்பது !!{enumerable} எனவும்,
$\cardless{\Nat}{A}$ என இருக்கும்போது, அப்போதும் அப்போது மட்டுமே,
$A$ என்பது !!{nonenumerable} எனவும் கூறலாம்.
இதனால் பின்வருவதை மீண்டும் கூற முடிகிறது:
\oliflabeldef{sfr:siz:nen-alt:thm:nonenum-pownat}{%
\olref[sfr][siz][nen-alt]{thm:nonenum-pownat}
என்பது
$\cardless{\Nat}{\Pow{\Nat}}$ என்ற கவனிப்பாகும்}{%
\olref[sfr][siz][nen]{thm:nonenum-pownat}
என்பது $\cardless{\PosInt}{\Pow{\PosInt}}$ என்ற கவனிப்பாகும்}.
உண்மையில், இந்தப் பிந்தைய கருத்து \emph{முழுப் பொதுமையுடையது}
என்று \citet{Cantor1892} நிறுவினார்:

% SEGMENT OLP-0036-S05
\begin{thm}[காண்டோர்]\ollabel{thm:cantor}
எந்தக் கணம் $A$ க்கும் $\cardless{A}{\Pow{A}}$.
\end{thm}

% SEGMENT OLP-0036-S06
\begin{proof}
$f(x) = \{x\}$ என்ற இணைப்பு,
!!a{injection} $f \colon A \to \Pow{A}$ ஆகும்.
ஏனெனில் $x \neq y$ எனில், உறுப்புசார் சமத்துவத்தின்படி
$\{x\} \neq \{y\}$ உம், எனவே $f(x) \neq f(y)$ உம் ஆகும்.
ஆகவே $\cardle{A}{\Pow{A}}$ எனப் பெற்றோம்.

\begin{editorial}
\olref[nen]{sec} இருந்தால் மெதுவான நிறுவலையும்,
இல்லையெனில் \olref[nen-alt]{sec} உடன் பொருந்தும் வேகமான
நிறுவலையும் தருகிறோம்.
\end{editorial}

\oliflabeldef{sfr:siz:nen:sec}{%
  !!a{surjective} சார்பு $g\colon A \to
  \Pow{A}$ இருக்கவே முடியாது, !!a{bijective} சார்பைப் பற்றிச்
  சொல்லவே வேண்டியதில்லை, என இப்போது காட்டுவோம்.
  இதனால் $\cardneq{A}{\Pow{A}}$.
  $g\colon A \to \Pow{A}$ எனக் கொள்வோம்.
  $g$ முழுச் சார்பாக இருப்பதால்,
  ஒவ்வொரு $x \in A$ உம் ஒரு உட்கணம் $g(x)
  \subseteq A$ க்கு அனுப்பப்படுகிறது.
  $g$ என்பது மேற்கோர்த்தல் பண்புடையதாக இருக்க முடியாது எனக் காட்டலாம்.
  இதற்கு, வரையறையின்படியே $g$ இன் வீச்சகத்தில் இருக்க முடியாத
  ஓர் உட்கணம் $\overline{A} \subseteq A$ ஐ வரையறுக்கிறோம்.
  பின்வருமாறு அமைக்க:
  \[
  \overline{A} = \Setabs{x \in A}{x \notin g(x)}.
  \]
  அனைத்து $x \in A$ க்கும் $g(x)$ வரையறுக்கப்பட்டுள்ளதால்,
  $\overline{A}$ என்பது தெளிவாக நன்கு வரையறுக்கப்பட்ட $A$ இன் உட்கணம்.
  ஆனால் அது $g$ இன் வீச்சகத்தில் இருக்க முடியாது.
  ஏதேனும் ஒரு $x \in A$ ஐக் கொண்டு,
  $\overline{A} \neq
  g(x)$ எனக் காட்டுவோம்.
  $x \in g(x)$ எனில், அது $x \notin g(x)$ என்பதை நிறைவேற்றாது.
  ஆகவே $\overline{A}$ இன் வரையறையின்படி,
  $x \notin
  \overline{A}$.
  $x \in \overline{A}$ எனில், அது $\overline{A}$ இன்
  வரையறுக்கும் பண்பை நிறைவேற்ற வேண்டும்;
  அதாவது $x \in A$, $x \notin
  g(x)$.
  % SOURCE CORRECTION OLSIZ-007; see translation/size-notes.tex.
  $x$ ஏதேனும் ஒன்றாக இருந்ததால், ஒவ்வொரு $x \in
  A$ க்கும், $x \in g(x)$ ஆக இருக்கும்போது, அப்போதும் அப்போது மட்டுமே,
  $x \notin \overline{A}$ என அமைகிறது;
  எனவே $g(x) \neq \overline{A}$.
  வேறு சொற்களில், $\overline{A}$ என்பது $g$ இன் வீச்சகத்தில்
  இருக்க முடியாது. இது $g$ மேற்கோர்த்தல் பண்புடையது என்ற
  முன்கோளுக்கு முரணாகும்.}{%
  $\cardneq{A}{\Pow{A}}$ எனக் காட்டுவது மீதமுள்ளது.
  முரண் பெறும் நோக்கில் $\cardeq{A}{\Pow{A}}$ எனக் கொள்வோம்.
  அதாவது ஏதோ ஓர் !!{bijection} $g \colon A \to \Pow{A}$ உள்ளது.
  இப்போது பின்வருவதைக் கருதுக:
  \[
    D = \Setabs{x \in A}{x \notin g(x)}
  \]
  $D \subseteq A$ என்பதையும், எனவே $D \in \Pow{A}$ என்பதையும் கவனிக்கவும்.
  $g$ என்பது !!a{bijection} என்பதால்,
  $g(y) = D$ என அமையும் ஏதோ ஒரு $y \in A$ உள்ளது.
  ஆனால் இப்போது:
  \[
    y \in g(y) \text{ அப்போதும் அப்போது மட்டுமே } y \in D
    \text{ அப்போதும் அப்போது மட்டுமே } y \notin g(y).
  \]
  இது ஒரு முரண்; ஆகவே $\cardneq{A}{\Pow{A}}$.}{}
\end{proof}

% SEGMENT OLP-0036-S07
\begin{explain}
\oliflabeldef{sfr:siz:nen:thm:nonenum-pownat}{%
  \olref{thm:cantor} இன் நிறுவலை,
  \olref[nen]{thm:nonenum-pownat} இன் நிறுவலுடன் ஒப்பிடுவது பயனுள்ளது.
  அங்கு $\PosInt$ இன் உட்கணங்களைக் கொண்ட எந்தப் பட்டியல்
  $Z_1$, $Z_2$, \dots க்கும்,
  அந்தப் பட்டியலில் இடம்பெறாது என்பது உறுதியான
  $\overline{Z}$ என்ற எண்களின் கணத்தைக் கட்டமைக்க முடியும் எனக் காட்டினோம்.
  ஒவ்வொரு $n \in
  \PosInt$ க்கும், $n \in Z_n$ ஆக இருக்கும்போது,
  அப்போதும் அப்போது மட்டுமே, $n \notin \overline{Z}$ என
  இருந்ததால் அது பட்டியலில் இல்லை என்பது உறுதியாயிற்று.
  இவ்வாறு, $Z_n$, $\overline{Z}$ ஆகியவற்றில் ஒன்றின்
  !!a{element} ஆகவும் மற்றொன்றின் உறுப்பு அல்லாததாகவும்
  ஏதாவது ஓர் எண் எப்போதும் உள்ளது.
  இங்கும் அதே கருத்தைப் பின்பற்றுகிறோம்;
  ஆனால் குறியீடுகள் $n$ இப்போது $\PosInt$ இன் உறுப்புகளாக இல்லாமல்,
  $A$ இன் !!{element}s ஆக உள்ளன.
  ஒவ்வொரு $x \in A$ க்கும் $\overline{A}$ என்பது $g(x)$ இலிருந்து
  வேறுபடுமாறு வரையறுக்கப்பட்டுள்ளது;
  ஏனெனில் $x \in g(x)$ ஆக இருக்கும்போது, அப்போதும் அப்போது மட்டுமே,
  $x \notin \overline{A}$.
  மீண்டும், $g(x)$, $\overline{A}$ ஆகியவற்றில் ஒன்றின்
  !!a{element} ஆகவும் மற்றொன்றின் உறுப்பு அல்லாததாகவும்
  $A$ இன் ஓர் !!a{element} எப்போதும் உள்ளது.
  எனவே $\overline{Z}$ என்பது $Z_1$, $Z_2$, \dots என்ற பட்டியலில்
  இடம்பெற முடியாததைப் போலவே,
  $\overline{A}$ என்பது $g$ இன் வீச்சகத்தில் இருக்க முடியாது.}{}

\oliflabeldef{sfr:siz:nen-alt:thm:nonenum-pownat}{%
  \olref{thm:cantor} இன் நிறுவலை,
  \olref[nen-alt]{thm:nonenum-pownat} இன் நிறுவலுடன் ஒப்பிடுவது பயனுள்ளது.
  அங்கு $\Nat$ இன் உட்கணங்களைக் கொண்ட எந்தப் பட்டியல்
  $N_0$, $N_1$, $N_2$, \dots க்கும்,
  அந்தப் பட்டியலில் இடம்பெறாது என்பது உறுதியான
  $D$ என்ற எண்களின் கணத்தைக் கட்டமைக்க முடியும் எனக் காட்டினோம்.
  ஒவ்வொரு $n \in \Nat$ க்கும்,
  $n \in N_n$ ஆக இருக்கும்போது, அப்போதும் அப்போது மட்டுமே,
  $n \notin
  D$ என இருந்ததால் அது பட்டியலில் இல்லை என்பது உறுதியாயிற்று.
  இங்கும் அதே கருத்தைப் பின்பற்றுகிறோம்;
  ஆனால் குறியீடுகள் $n$ இப்போது $\Nat$ இன் உறுப்புகளாக இல்லாமல்,
  $A$ இன் !!{element}s ஆக உள்ளன.
  ஒவ்வொரு $x
  \in A$ க்கும் $D$ என்பது $g(x)$ இலிருந்து வேறுபடுமாறு
  வரையறுக்கப்பட்டுள்ளது;
  ஏனெனில் $x \in g(x)$ ஆக இருக்கும்போது, அப்போதும் அப்போது மட்டுமே,
  $x \notin D$.}{}

இந்த நிறுவலை ரஸ்ஸலின் முரண்பாட்டின் நிறுவலான
\olref[sfr][set][rus]{thm:russells-paradox} உடன் ஒப்பிடுவதும் பயனுள்ளது.
உண்மையில், காண்டோரின் தேற்றமே ரஸ்ஸலின் முரண்பாட்டிற்கு
உந்துதலாக அமைந்தது.
\end{explain}

% SEGMENT OLP-0036-S08
\begin{prob}
எந்தக் கணம் $A$ க்கும் !!a{injection}
$g\colon \Pow{A} \to
A$ இருக்க முடியாது எனக் காட்டுக.
குறிப்பு: $g\colon \Pow{A} \to A$ என்பது !!{injective}
பண்புடையது எனக் கொள்க.
$D = \Setabs{g(B)}{B \subseteq A \text{ மற்றும்
} g(B) \notin B}$ என்பதைக் கருதுக.
$x = g(D)$ என அமைக்க.
$g$ என்பது !!{injective} பண்புடையது என்ற உண்மையைப் பயன்படுத்தி
ஒரு முரணைப் பெறுக.
\end{prob}

\end{document}
